% Op-Amp Circuits
% Demonstrates: inverting, non-inverting, and differential amplifier configurations.
% Compile with: pdflatex main.tex

\documentclass[border=10pt]{standalone}
\usepackage[american]{circuitikz}
\usepackage{amsmath}

\begin{document}
\begin{circuitikz}[scale=0.85, transform shape]

% =============================================================================
% Circuit 1: Inverting Amplifier
% =============================================================================
\begin{scope}[shift={(0, 0)}]
    \node[font=\bfseries\large] at (3, 4) {Inverting Amplifier};

    % Op-amp
    \node[op amp, noinv input up] (opamp) at (3, 1.5) {};

    % Input resistor R1
    \draw (0, 2) node[left] {$V_{\text{in}}$}
        to[R, l=$R_1$] (opamp.-)  ;

    % Feedback resistor Rf
    \draw (opamp.-) -- ++(0, 1.5) coordinate (fb)
        to[R, l=$R_f$] (fb -| opamp.out) -- (opamp.out);

    % Non-inverting input to ground
    \draw (opamp.+) -- ++(0, -0.5) node[ground] {};

    % Output
    \draw (opamp.out) -- ++(1.5, 0) node[right] {$V_{\text{out}} = -\dfrac{R_f}{R_1}\,V_{\text{in}}$};
\end{scope}

% =============================================================================
% Circuit 2: Non-Inverting Amplifier
% =============================================================================
\begin{scope}[shift={(10, 0)}]
    \node[font=\bfseries\large] at (3, 4) {Non-Inverting Amplifier};

    % Op-amp
    \node[op amp, noinv input up] (opamp) at (3, 1.5) {};

    % Input to non-inverting
    \draw (opamp.+) -- ++(-1.5, 0) node[left] {$V_{\text{in}}$};

    % Feedback network
    \draw (opamp.-) -- ++(0, -1) coordinate (gnd)
        to[R, l_=$R_1$] (gnd |- 0,-1) node[ground] {};

    \draw (opamp.-) -- ++(0, 1.5) coordinate (fb)
        to[R, l=$R_f$] (fb -| opamp.out) -- (opamp.out);

    % Output
    \draw (opamp.out) -- ++(1.5, 0)
        node[right] {$V_{\text{out}} = \left(1 + \dfrac{R_f}{R_1}\right)V_{\text{in}}$};
\end{scope}

% =============================================================================
% Circuit 3: Voltage Follower (Buffer)
% =============================================================================
\begin{scope}[shift={(0, -7)}]
    \node[font=\bfseries\large] at (3, 4) {Voltage Follower};

    % Op-amp
    \node[op amp, noinv input up] (opamp) at (3, 1.5) {};

    % Input
    \draw (opamp.+) -- ++(-1.5, 0) node[left] {$V_{\text{in}}$};

    % Direct feedback (output to inverting)
    \draw (opamp.out) -- ++(0.5, 0) coordinate (out)
        -- ++(0, -2) -| (opamp.-);

    % Output
    \draw (out) -- ++(1, 0) node[right] {$V_{\text{out}} = V_{\text{in}}$};
\end{scope}

% =============================================================================
% Circuit 4: Summing Amplifier
% =============================================================================
\begin{scope}[shift={(10, -7)}]
    \node[font=\bfseries\large] at (3, 4.5) {Summing Amplifier};

    % Op-amp
    \node[op amp, noinv input up] (opamp) at (4, 1.5) {};

    % Input resistors
    \draw (0, 3) node[left] {$V_1$}
        to[R, l=$R_1$] (2.5, 3) -- (2.5, 2);
    \draw (0, 2) node[left] {$V_2$}
        to[R, l=$R_2$] (2.5, 2);
    \draw (0, 1) node[left] {$V_3$}
        to[R, l=$R_3$] (2.5, 1) -- (2.5, 2);

    % Connect to inverting input
    \draw (2.5, 2) -- (opamp.-);

    % Feedback resistor
    \draw (opamp.-) ++(0, 0) -- ++(0, 1.5) coordinate (fb)
        to[R, l=$R_f$] (fb -| opamp.out) -- (opamp.out);

    % Non-inverting to ground
    \draw (opamp.+) -- ++(0, -0.5) node[ground] {};

    % Output
    \draw (opamp.out) -- ++(0.8, 0)
        node[right, font=\small] {$V_{\text{out}} = -R_f\!\left(\dfrac{V_1}{R_1}+\dfrac{V_2}{R_2}+\dfrac{V_3}{R_3}\right)$};
\end{scope}

\end{circuitikz}
\end{document}
